%==============================================================
% CrossVul-DG: Class-Conditional Domain Generalization for
% Cross-Project Software Vulnerability Detection
% IEEE conference format, double-blind.
%
% NOTE: Every quantitative value below is an illustrative
% placeholder until scripts/run_full_sweep.sh has been run on
% MegaVul. Replace the em-dashes / placeholder numbers with
% measured results before submission.
%==============================================================
\documentclass[conference]{IEEEtran}
\IEEEoverridecommandlockouts

\usepackage{cite}
\usepackage{amsmath,amssymb,amsfonts}
\usepackage{algorithm}
\usepackage{algpseudocode}
\usepackage{graphicx}
\usepackage{booktabs}
\usepackage{url}
\usepackage{xcolor}
\usepackage[colorlinks=true,allcolors=blue]{hyperref}

% Lightweight "placeholder/notice" box used to keep the draft honest.
\newcommand{\notice}[1]{%
  \par\vspace{2pt}\noindent
  \fcolorbox{red!60}{red!4}{\parbox{0.97\columnwidth}{\footnotesize
  \textbf{Placeholder notice.} #1}}\par\vspace{3pt}}

\begin{document}

\title{CrossVul-DG: Class-Conditional Domain Generalization\\
for Cross-Project Software Vulnerability Detection}

\author{\IEEEauthorblockN{Anonymous Author(s)}
\IEEEauthorblockA{Affiliation withheld for double-blind review\\
Email: anonymous@example.com}}

\maketitle

%==============================================================
\begin{abstract}
Deep-learning models for software vulnerability detection (SVD)
degrade sharply on projects unseen during training, because of
domain shift in coding conventions, naming, and architecture.
Domain Adaptation (DA) needs target-project data at training
time, which is unrealistic for proactive deployment; recent
Domain Generalization (DG) for SVD takes the \emph{project} as
the only unit of invariance and is \emph{class-blind}, ignoring
the severe class imbalance of real vulnerability data. We argue
cross-project SVD is a \emph{coupled} domain-shift and
class-imbalance problem, and we propose \textbf{CC-DRO}
(Class-Conditional Distributionally-Robust Optimization): it
forms distributional groups over \emph{(project $\times$ class)}
cells and adds a cross-project \emph{positive-class} risk-invariance
penalty, so a model's vulnerability sensitivity is robust
\emph{and} consistent across unseen projects. We study CC-DRO
against ERM, IRM, GroupDRO, two DA baselines (which access
unlabelled target data), and a pseudo-domain variant, under a
leave-one-project-out (LOPO) protocol on MegaVul with a
chronological 8:1:1 split, a shared imbalanced-data protocol
(class-weighted loss, validation-tuned threshold, PR-AUC as the
primary metric), and three encoder backbones. We expect CC-DRO
to improve cross-project PR-AUC and, in particular, to narrow
the cross-project \emph{vulnerability-recall gap} relative to
class-blind DG, with only a marginal seen-project regression.
The figures in this manuscript are illustrative placeholders
pending the full run; we release code and pre-processed splits
for reproducibility.
\end{abstract}

\begin{IEEEkeywords}
software vulnerability detection, domain generalization,
class imbalance, distributionally robust optimization, IRM,
MegaVul, cross-project, CodeBERT
\end{IEEEkeywords}

%==============================================================
\section{Introduction}
Software vulnerabilities cause large-scale economic and societal
harm, and automated detection is an intensive research focus,
with deep-learning detectors achieving strong in-distribution
F1~\cite{vuldeepecker,devign,linevul,vulberta}. Yet when these
models are trained on one codebase and applied to
another, performance often drops
substantially~\cite{reveal,steenhoek}. This cross-project gap is
a fundamental barrier to deployment: a security team rarely has
labelled vulnerabilities in the target codebase before scanning it.

\notice{This is a study-protocol draft. Every quantitative value
in the abstract and in Section~\ref{sec:results} is an
illustrative placeholder produced only to exercise the released
analysis pipeline; the numbers are \emph{not} measured results on
MegaVul. They must be replaced with values from the experiments
in Section~\ref{sec:setup} before this manuscript is submitted.}

A natural response is Domain Adaptation (DA), aligning the model
using unlabelled or partially labelled target
data~\cite{dann,coral}. DA presupposes target data at training
time, which is unrealistic for proactive deployment (e.g., a
vendor shipping a detector to unknown customer codebases). Domain
Generalization (DG) instead learns representations that
generalise to \emph{any} unseen domain using only source data at
training time~\cite{muandet,irm,dro}. DG is mature in computer
vision; in SVD it has begun to appear only recently, e.g.\ the
cluster-based ensemble method CEVD~\cite{cevd}.

\noindent\textbf{The coupled problem.}
Cross-project SVD is \emph{not} a pure domain-shift problem. Real
vulnerability datasets are highly imbalanced --- in MegaVul
roughly one function in twenty is vulnerable. Off-the-shelf DG
objectives (IRM, GroupDRO) take the project as the unit of
invariance and are \emph{class-blind}: under $\sim$1:19 imbalance,
the worst-case \emph{per-project} risk can be minimised by
excelling on the dominant safe class, so the rare-but-critical
vulnerability signal is under-served, and nothing in the objective
makes vulnerability \emph{recall} consistent across projects. A
detector can thus look strong on aggregate while its ability to
flag real vulnerabilities swings widely from project to project.

\noindent\textbf{This paper.}
We present CrossVul-DG, an empirical study of
invariance/robustness-based DG for cross-project SVD on
MegaVul~\cite{megavul}, and we introduce \textbf{CC-DRO}, a
class-conditional robust objective for the coupled problem.
CC-DRO (i) replaces project-level groups with (project
$\times$ class) cells so the worst case explicitly includes the
\emph{(project, vulnerable)} cells, and (ii) adds a penalty equal
to the variance, across projects, of the positive-class risk,
forcing vulnerability sensitivity to be invariant across projects.
All methods share a single imbalanced-data protocol so the
comparison is fair, and we test generality across three code
encoders (CodeBERT~\cite{codebert},
GraphCodeBERT~\cite{graphcodebert}, UniXcoder~\cite{unixcoder}).

\noindent\textbf{Research questions.}
\begin{itemize}
\item \textbf{RQ1.} Does DG (IRM, GroupDRO, or CC-DRO) improve
cross-project PR-AUC over standard ERM fine-tuning?
\item \textbf{RQ2.} How does DG compare with DA methods that have
access to (unlabelled) target data?
\item \textbf{RQ3.} \emph{(Core)} Does class-conditional
invariance (CC-DRO) reduce the cross-project
\emph{vulnerability-recall gap} relative to class-blind DG?
\item \textbf{RQ4.} What is the trade-off between cross-project
gain and seen-project performance?
\item \textbf{RQ5.} Does the CC-DRO gain transfer across encoder
backbones?
\item \textbf{RQ6.} \emph{(Optional)} How do prompted code-LLMs
perform as zero-/few-shot cross-project detectors?
\end{itemize}

\noindent\textbf{Contributions.}
\begin{enumerate}
\item We frame cross-project SVD as a \emph{coupled} domain-shift
and class-imbalance problem and propose \textbf{CC-DRO}, to our
knowledge the first class-conditional robust DG objective for SVD,
with an explicit cross-project positive-class invariance penalty.
\item We provide a reproducible LOPO protocol on MegaVul with a
shared imbalanced-data treatment (class-weighted loss,
validation-tuned decision threshold, PR-AUC primary metric, and a
leakage-free DA target split), and statistical significance
testing.
\item We adapt IRM, GroupDRO and a pseudo-domain extension to the
same protocol as baselines, and evaluate across three code
encoders.
\item We release pre-processed splits, training pipelines, and
evaluation scripts.\footnote{\url{https://github.com/anonymous/crossvul-dg}
(anonymised for review).} (Quantitative gains in this draft are
placeholders to be confirmed by the released pipeline.)
\end{enumerate}

%==============================================================
\section{Background and Related Work}
\subsection{Deep Learning for Vulnerability Detection}
Early sequence detectors such as VulDeePecker~\cite{vuldeepecker}
used LSTMs over token sequences. Graph methods, notably
Devign~\cite{devign} and the ReVeal study~\cite{reveal}, encoded
code property graphs with GNNs. Transformer encoders ---
CodeBERT~\cite{codebert}, GraphCodeBERT~\cite{graphcodebert},
UniXcoder~\cite{unixcoder} --- are now the de facto backbone, with
LineVul~\cite{linevul} a strong line-level baseline and
VulBERTa~\cite{vulberta} a simplified recipe. LLM
studies~\cite{vulbench} report that prompted detectors remain
expensive and unstable across projects.

\subsection{Cross-Project Vulnerability Detection}
Chakraborty et al.~\cite{reveal} documented the severity of
cross-project drops; Steenhoek et al.~\cite{steenhoek} confirmed
them under a controlled study. Two recent lines target
generalization directly. CEVD~\cite{cevd} is a DG method that
clusters code embeddings into pseudo-domains and trains an
ensemble of per-domain experts fused by a domain classifier, with
no target data. VulGate~\cite{vulgate} is an orthogonal,
data-centric effort: a large de-duplicated dataset (236,663
functions, with a 500-sample expert-verified cross-project test
set) showing data quality and model choice strongly affect
generalization. CrossVul-DG differs from CEVD by applying explicit
invariance/robustness objectives over the natural project
environments and, with CC-DRO, by conditioning the robust
objective on \emph{class}; it is complementary to VulGate, which
improves the data rather than the objective.

\subsection{Domain Generalization and Class Imbalance}
DG learns predictors that generalise to unseen
domains~\cite{muandet}. IRM~\cite{irm} penalises representations
on which the optimal classifier differs across environments;
GroupDRO~\cite{dro} minimises the worst-case group risk;
gradient-matching (Fish~\cite{fish}) and Mixup-style augmentations
extend the family. Surveys~\cite{dgsurvey1,dgsurvey2} note DG's
maturity in vision/NLP but limited adoption in software
engineering. Crucially, all of these define groups by domain
\emph{only}. CC-DRO defines groups by domain $\times$ class and
adds a positive-class invariance term --- a class-conditional
robust objective tailored to imbalanced SVD.

%==============================================================
\section{The CrossVul-DG Framework}
\subsection{Problem Formulation}
Let $\mathcal{P}=\{p_1,\dots,p_K\}$ be the source projects. Each
$p_k$ provides $D_k=\{(x_i^{(k)},y_i^{(k)})\}$ with function $x_i$
and label $y_i\in\{0,1\}$ (vulnerable or not). The positive rate
$\pi_k=\Pr[y=1\mid p_k]$ is small ($\approx 0.05$). We seek
$f_\theta:\mathcal{X}\to\{0,1\}$ that generalises to an unseen
project $p_\ast\notin\mathcal{P}$ using only
$\bigcup_k D_k$ at training time (no target access), and whose
\emph{positive-class} performance is consistent across projects.

\subsection{Backbone Encoder}
We use a fine-tuned code encoder $g_\phi$ (CodeBERT by default;
GraphCodeBERT and UniXcoder for generality) producing a 768-d
first-token representation, followed by a two-layer MLP $h_\psi$
with a single logit output. We write $f_\theta=h_\psi\circ g_\phi$,
$\theta=(\phi,\psi)$.

\subsection{Imbalanced-Setting Protocol (shared by all methods)}
\label{sec:imbal}
Because we keep MegaVul's natural $\approx$1:19 ratio, every method
shares the following, so differences reflect the objective, not
the imbalance treatment:
\begin{itemize}
\item \textbf{Class-weighted loss.} Per-sample
$\mathrm{BCE}$ with $\mathrm{pos\_weight}=N_{\text{neg}}/N_{\text{pos}}$
on the training fold (capped), so optimisation does not collapse to
the safe class.
\item \textbf{Model selection on validation PR-AUC} (threshold-free;
validation-F1 at a fixed $0.5$ threshold is unreliable under
imbalance).
\item \textbf{Validation-tuned decision threshold} (max-F1),
applied unchanged at test --- no test-set peeking.
\item \textbf{Primary metric PR-AUC}, with F1 at the tuned
threshold, recall at precision $\ge 0.5$ (R@P50), and
\emph{worst-project F1}.
\item \textbf{Leakage-free DA target.} DA baselines use the
chronologically \emph{earlier} portion of the held-out project as
the unlabelled target; \emph{all} methods are scored on the
\emph{later} portion (identical test set).
\end{itemize}

\subsection{ERM, IRM, GroupDRO (baselines)}
Let $R_k(\theta)=\mathbb{E}_{D_k}[\,w_y\,\ell(f_\theta(x),y)\,]$ be
the class-weighted risk on project $k$ ($w_y$ the pos\_weight on
positives). ERM minimises $\frac1K\sum_k R_k$. IRM~\cite{irm} adds
\begin{equation}
\mathcal{L}_{\text{IRM}}=\sum_{k=1}^{K}R_k(\theta)
+\lambda\sum_{k=1}^{K}\big\|\nabla_{w|w=1}R_k(w\!\cdot\! f_\theta)\big\|^2,
\end{equation}
ramping $\lambda$ after a warm-up. GroupDRO~\cite{dro} minimises
the worst-case project risk
$\min_\theta\max_{q\in\Delta_K}\sum_k q_k R_k(\theta)$ via
exponentiated-gradient updates
$q_k\leftarrow q_k\exp(\eta_q R_k);\, q\leftarrow q/\sum_{k'}q_{k'}$.

\subsection{CC-DRO: Class-Conditional DRO (proposed)}
\label{sec:ccdro}
CC-DRO changes the group structure and adds an invariance term.
\textbf{(1) Class-conditional groups.} Groups are cells
$g=(p,c)$ for project $p$ and class $c\in\{0,1\}$, giving $2K$
cells; GroupDRO is run over the cells, so the worst case
explicitly includes the $(p,1)$ \emph{vulnerable} cells.
\textbf{(2) Positive-class invariance.} With $R_{p,\text{pos}}$ the
mean weighted loss on project $p$'s positive samples, we penalise
their variance across projects. The per-step objective is
\begin{equation}
\label{eq:ccdro}
\mathcal{L}_{\text{CC-DRO}}=\sum_{g} q_g R_g(\theta)
\;+\;\lambda_{\text{inv}}\;\operatorname{Var}_{p}\!\big(R_{p,\text{pos}}\big),
\end{equation}
where $g$ ranges over the cells present in the step and $q$ are
exp-gradient weights over cells. Minimising the worst-case cell
risk concentrates capacity on the hardest \emph{(project,
vulnerable)} cells; minimising the positive-risk variance makes
vulnerability recall consistent across projects. Cells absent from
a mini-batch are omitted that step (rare-positive projects may
yield no positives in a small batch); the penalty is applied only
when $\ge 2$ projects contribute positives, and $\lambda_{\text{inv}}$
is ramped after a warm-up. Algorithm~\ref{alg:ccdro} gives one
optimisation step.

\begin{algorithm}[t]
\caption{CC-DRO optimisation step}
\label{alg:ccdro}
\begin{algorithmic}[1]
\Require per-project batches $\{B_k\}_{k=1}^{K}$; cell weights
$q$; pos\_weight $w$; penalty weight $\lambda_{\text{inv}}$
\State $\mathcal{C}\gets\{\}$ \Comment{cell $\to$ mean loss};
$\;\mathcal{R}^{+}\gets[\,]$
\For{each project $k$ with batch $B_k$}
  \State $\ell_i \gets w^{\,y_i}\,\mathrm{BCE}(f_\theta(x_i),y_i)$
  for $(x_i,y_i)\in B_k$
  \For{$c\in\{0,1\}$ present in $B_k$}
     \State $\mathcal{C}[(k,c)] \gets \mathrm{mean}\{\ell_i : y_i=c\}$
  \EndFor
  \If{$(k,1)\in\mathcal{C}$}
     \State append $\mathcal{C}[(k,1)]$ to $\mathcal{R}^{+}$
  \EndIf
\EndFor
\State $q_g \gets q_g\exp(\eta_q\,\mathcal{C}[g])$ for $g\in\mathcal{C}$;\;
$q\gets q/\sum q$ \Comment{DRO update}
\State $L \gets \sum_{g\in\mathcal{C}} q_g\,\mathcal{C}[g]$
\If{$|\mathcal{R}^{+}|\ge 2$}
  \State $L \gets L + \lambda_{\text{inv}}\operatorname{Var}(\mathcal{R}^{+})$
\EndIf
\State back-propagate $L$; update $\theta$
\end{algorithmic}
\end{algorithm}

\subsection{Pseudo-Domain Clustering (long-tail extension)}
MegaVul has a long tail: the top-10 projects cover most samples
while $\sim$980 are small. We optionally cluster small projects in
embedding space (mean function embedding per project) into $C$
pseudo-domains with $k$-means and add them as extra environments
for GroupDRO (PD+DRO). Clustering code embeddings to expose latent
domains is also central to CEVD~\cite{cevd}; here it is an
auxiliary source of environments for a robust objective rather
than the basis for an ensemble. Algorithm~\ref{alg:pd} summarises
the procedure.

\begin{algorithm}[t]
\caption{Pseudo-domain construction (PD+DRO)}
\label{alg:pd}
\begin{algorithmic}[1]
\Require top-$K$ projects $\mathcal{P}_{\text{top}}$; tail projects
$\mathcal{P}_{\text{tail}}$; cluster count $C$
\State compute mean embedding $\bar z_p$ for each
$p\in\mathcal{P}_{\text{tail}}$ via the encoder
\State cluster $\{\bar z_p\}$ into $C$ groups with $k$-means
\State environments
$\mathcal{E}\gets\mathcal{P}_{\text{top}}\cup\{c_1,\dots,c_C\}$
\State train $f_\theta$ by minimising GroupDRO over $\mathcal{E}$
\State \Return $f_\theta$
\end{algorithmic}
\end{algorithm}

%==============================================================
\section{Experimental Setup}
\label{sec:setup}
\subsection{Dataset}
We use MegaVul~\cite{megavul}, a large real-world C/C++ dataset:
17,380 vulnerabilities across 169 CWE types (2006--2023) from 992
repositories; the function-level release has $\approx$17,000
vulnerable and $\approx$320,000 non-vulnerable functions with CVE
identifiers and commit timestamps. We select the top-10 projects
by function count as LOPO environments (Table~\ref{tab:projects});
the exact post-deduplication share is produced by the
preprocessing script.

\begin{table}[t]
\centering
\caption{Top-10 MegaVul projects used as LOPO environments.
\emph{Placeholder: exact post-deduplication counts are produced by
\texttt{prepare\_data.py} and will replace the dashes.}}
\label{tab:projects}
\begin{tabular}{lccc}
\toprule
Project & \#Func. & \#Vuln. & Vuln.\ \% \\
\midrule
Linux & --- & --- & --- \\
Chromium & --- & --- & --- \\
FFmpeg & --- & --- & --- \\
QEMU & --- & --- & --- \\
PHP-src & --- & --- & --- \\
ImageMagick & --- & --- & --- \\
OpenSSL & --- & --- & --- \\
TensorFlow & --- & --- & --- \\
Wireshark & --- & --- & --- \\
systemd & --- & --- & --- \\
\bottomrule
\end{tabular}
\end{table}

\subsection{Splits and Protocol}
A chronological 8:1:1 train/val/test split by commit date over the
pooled selected samples (cut at the 0.8 and 0.9 time quantiles),
rather than by year. For cross-project evaluation we use LOPO: each
top-10 project is held out in turn; the held-out project is split
chronologically into an earlier DA-target portion and a later test
portion (Section~\ref{sec:imbal}). We report mean $\pm$ std across
seeds (three for the main table; two for the backbone study) and
test significance with a paired Wilcoxon signed-rank test over the
10 folds.

\subsection{Baselines and Methods}
\textbf{ERM-FT} (fine-tuning, ignores project labels);
\textbf{DA-DANN}~\cite{dann} and \textbf{DA-CORAL}~\cite{coral}
(use the earlier portion of the held-out project as unlabelled
target --- a reference, not strict DG); \textbf{CrossVul-DG (IRM)},
\textbf{(DRO)}; \textbf{CC-DRO} (proposed); and
\textbf{(PD+DRO)} (pseudo-domain + GroupDRO).

\subsection{Implementation Details}
Encoder \texttt{microsoft/codebert-base} (and
\texttt{unixcoder-base}, \texttt{graphcodebert-base} for Tier~2);
sequence length 512; batch size 32 (per-environment for DG);
AdamW, lr $2\times10^{-5}$, weight decay $10^{-2}$; 10 epochs,
early stopping on validation PR-AUC; IRM $\lambda_{\max}{=}10$,
warm-up 1 epoch; GroupDRO/CC-DRO $\eta_q{=}0.01$; CC-DRO
$\lambda_{\text{inv}}{=}1$ (ablated in $\{0,0.1,1,10\}$), warm-up 1
epoch; pseudo-domain $C{=}5$. All runs fit a single 24~GB GPU.

\subsection{Metrics}
PR-AUC (primary), F1 at the validation-tuned threshold, recall at
precision $\ge0.5$ (R@P50), ROC-AUC, and worst-project F1.
Following~\cite{reveal} we also report the per-project F1 range.

%==============================================================
\section{Expected Results (Illustrative)}
\label{sec:results}
\notice{Every number in this section is a placeholder obtained by
running the released pipeline on tiny synthetic data, included only
to show the reporting format and verify the analysis scripts
end-to-end. These are \emph{not} measured results on MegaVul.
Replace them with the output of \texttt{run\_full\_sweep.sh} before
drawing any conclusion or submitting.}

Tables~\ref{tab:main}--\ref{tab:backbone} illustrate the intended
presentation and the hypothesised direction of effects
(CC-DRO $\ge$ class-blind DG $\ge$ DA $\ge$ ERM on cross-project
PR-AUC, and a smaller vulnerability-recall gap for CC-DRO). They do
not constitute evidence until the experiments are executed.

\begin{table}[t]
\centering
\caption{Leave-one-project-out results on MegaVul (CodeBERT),
mean over 10 folds (each averaged over seeds). PR-AUC is primary;
worst-F1 is the robustness number. $^\dagger$ uses unlabelled
target data. \emph{Illustrative placeholders --- not measured.}}
\label{tab:main}
\begin{tabular}{lcccc}
\toprule
Method & PR-AUC & F1$_{\text{thr}}$ & R@P50 & worst-F1 \\
\midrule
ERM-FT             & --- & --- & --- & --- \\
DA-CORAL$^\dagger$ & --- & --- & --- & --- \\
DA-DANN$^\dagger$  & --- & --- & --- & --- \\
IRM                & --- & --- & --- & --- \\
GroupDRO           & --- & --- & --- & --- \\
\textbf{CC-DRO (ours)} & --- & --- & --- & --- \\
PD+DRO             & --- & --- & --- & --- \\
\bottomrule
\end{tabular}
\end{table}

\begin{table}[t]
\centering
\caption{RQ3: cross-project vulnerability-recall gap (max$-$min
per-project recall at the tuned threshold; lower is better) and
mean vulnerability recall. \emph{Placeholders --- not measured.}}
\label{tab:gap}
\begin{tabular}{lcc}
\toprule
Method & Recall gap $\downarrow$ & Mean recall $\uparrow$ \\
\midrule
IRM      & --- & --- \\
GroupDRO & --- & --- \\
\textbf{CC-DRO (ours)} & --- & --- \\
\bottomrule
\end{tabular}
\end{table}

\begin{table}[t]
\centering
\caption{Seen vs.\ unseen project F1 (averaged over the 10 LOPO
folds). \emph{Placeholders --- not measured.}}
\label{tab:seen}
\begin{tabular}{lcc}
\toprule
Method & Seen F1 & Unseen F1 \\
\midrule
ERM-FT & --- & --- \\
\textbf{CC-DRO (ours)} & --- & --- \\
$\Delta$ & --- & --- \\
\bottomrule
\end{tabular}
\end{table}

\begin{table}[t]
\centering
\caption{RQ5: backbone generality (PR-AUC, mean over 10 folds).
\emph{Placeholders --- not measured.}}
\label{tab:backbone}
\begin{tabular}{lccc}
\toprule
Backbone & ERM-FT & GroupDRO & \textbf{CC-DRO} \\
\midrule
CodeBERT      & --- & --- & --- \\
GraphCodeBERT & --- & --- & --- \\
UniXcoder     & --- & --- & --- \\
\bottomrule
\end{tabular}
\end{table}

\begin{figure}[t]
\centering
\fbox{\parbox{0.92\columnwidth}{\centering\footnotesize
\vspace{18pt} \emph{Fig.~1 placeholder:} ablation of CC-DRO
$\lambda_{\text{inv}}\!\in\!\{0,0.1,1,10\}$ (PR-AUC and recall gap)
and pseudo-domain count $C\!\in\!\{0,2,5,10,20\}$. \vspace{18pt}}}
\caption{Planned ablations; curves are schematic until measured.}
\label{fig:abl}
\end{figure}

\subsection{RQ1: DG vs.\ ERM}
Table~\ref{tab:main} shows the intended LOPO summary, with
placeholders consistent with the hypothesis that robust DG (and
CC-DRO in particular) improves cross-project PR-AUC over ERM.

\subsection{RQ2: DG vs.\ DA}
The key question is whether DG can match DA despite seeing no
target labels. The placeholders encode the hypothesised outcome
(DG $\ge$ DA); a plausible mechanism, if confirmed, is that with
many natural source domains, invariance objectives extract more
transferable signal than adversarial alignment to a single target.

\subsection{RQ3: Vulnerability-recall gap (core)}
Table~\ref{tab:gap} is the central evaluation: we expect CC-DRO to
reduce the per-project recall gap relative to class-blind IRM and
GroupDRO, by construction of its class-conditional groups and
positive-class invariance penalty (Eq.~\ref{eq:ccdro}).

\subsection{RQ4: Seen vs.\ unseen trade-off}
Table~\ref{tab:seen} reports the hypothesised small seen-project
regression in exchange for a larger unseen-project gain.

\subsection{RQ5: Backbone generality}
Table~\ref{tab:backbone} tests whether the CC-DRO gain holds, in
direction, across CodeBERT, GraphCodeBERT, and UniXcoder.

%==============================================================
\section{Discussion}
\noindent\textbf{Why might class-conditioning help?}
Under heavy imbalance, a domain-only robust objective can satisfy
its worst-case constraint through the majority class. Conditioning
groups on class forces the optimiser to confront the
\emph{(project, vulnerable)} cells directly, and the positive-risk
variance term discourages solutions that detect vulnerabilities
well on some projects and poorly on others. Whether this
materialises on MegaVul is an empirical question.

\noindent\textbf{Practical implications (if confirmed).}
Tools shipped to many customer codebases face exactly the DG
setting. If the hypotheses hold, practitioners should (i) gather as
many distinct source projects as possible, (ii) prefer
class-conditional robust objectives under imbalance, and (iii)
accept a small seen-project regression for substantial
unseen-project gains in vulnerability recall.

\noindent\textbf{Threats to validity.}
\emph{Internal:} IRM/DRO are hyperparameter-sensitive; we report
mean $\pm$ std and provide CC-DRO ablations. \emph{External:}
findings are bound to C/C++ MegaVul; cross-language generalisation
is future work. \emph{Construct:} MegaVul labels inherit upstream
noise; the leakage-free DA split reduces the per-fold test size, a
trade-off we report. \emph{Conclusion:} until the protocol is run,
the numbers are illustrative and support no conclusions.

%==============================================================
\section{Conclusion}
We recast cross-project SVD as a coupled domain-shift and
class-imbalance problem and proposed \textbf{CC-DRO}, a
class-conditional robust DG objective with a cross-project
positive-class invariance penalty. Within a reproducible LOPO
protocol on MegaVul --- shared imbalanced-data treatment, PR-AUC
primary metric, leakage-free DA comparison, and three backbones ---
the study will test whether class-conditioning narrows the
cross-project vulnerability-recall gap left by class-blind DG, and
whether DG can rival DA without target access. It complements
cluster-based DG (CEVD~\cite{cevd}) and the data-centric VulGate
effort~\cite{vulgate}. Future work includes cross-language
settings, causal feature discovery, and combining DG with
retrieval-augmented LLM detectors.

%==============================================================
\section*{Data and Code Availability}
Pre-processed splits, training pipelines, and evaluation scripts
are available at \url{https://github.com/anonymous/crossvul-dg}
(anonymised for review).

%==============================================================
\begin{thebibliography}{00}
\bibitem{cevd} J. Liu, ``CEVD: Cluster-based ensemble learning for cross-project vulnerability detection,'' \emph{Future Internet}, vol.~18, no.~2, p.~85, 2026.
\bibitem{vuldeepecker} Z. Li, D. Zou, S. Xu, X. Ou, H. Jin, S. Wang, Z. Deng, and Y. Zhong, ``VulDeePecker: A deep learning-based system for vulnerability detection,'' in \emph{Proc. NDSS}, 2018.
\bibitem{devign} Y. Zhou, S. Liu, J. Siow, X. Du, and Y. Liu, ``Devign: Effective vulnerability identification by learning comprehensive program semantics via graph neural networks,'' in \emph{NeurIPS}, 2019.
\bibitem{linevul} M. Fu and C. Tantithamthavorn, ``LineVul: A transformer-based line-level vulnerability prediction,'' in \emph{Proc. MSR}, 2022, pp. 608--620.
\bibitem{vulberta} H. Hanif and S. Maffeis, ``VulBERTa: Simplified source code pre-training for vulnerability detection,'' in \emph{Proc. IJCNN}, 2022.
\bibitem{reveal} S. Chakraborty, R. Krishna, Y. Ding, and B. Ray, ``Deep learning based vulnerability detection: Are we there yet?'' \emph{IEEE Trans. Softw. Eng.}, vol.~48, no.~9, pp.~3280--3296, 2022.
\bibitem{steenhoek} B. Steenhoek, M. M. Rahman, R. Jiles, and W. Le, ``An empirical study of deep learning models for vulnerability detection,'' in \emph{Proc. ICSE}, 2023, pp.~2237--2248.
\bibitem{dann} Y. Ganin and V. Lempitsky, ``Unsupervised domain adaptation by backpropagation,'' in \emph{ICML}, 2015, pp.~1180--1189.
\bibitem{coral} B. Sun and K. Saenko, ``Deep CORAL: Correlation alignment for deep domain adaptation,'' in \emph{ECCV Workshops}, 2016, pp.~443--450.
\bibitem{irm} M. Arjovsky, L. Bottou, I. Gulrajani, and D. Lopez-Paz, ``Invariant risk minimization,'' \emph{arXiv:1907.02893}, 2019.
\bibitem{dro} S. Sagawa, P. W. Koh, T. B. Hashimoto, and P. Liang, ``Distributionally robust neural networks for group shifts,'' in \emph{ICLR}, 2020.
\bibitem{muandet} K. Muandet, D. Balduzzi, and B. Sch\"olkopf, ``Domain generalization via invariant feature representation,'' in \emph{ICML}, 2013, pp.~10--18.
\bibitem{megavul} C. Ni, L. Shen, X. Yang, Y. Zhu, and S. Wang, ``MegaVul: A C/C++ vulnerability dataset with comprehensive code representations,'' in \emph{Proc. MSR}. IEEE, 2024, pp.~738--742.
\bibitem{codebert} Z. Feng, D. Guo, D. Tang, N. Duan, X. Feng, M. Gong, L. Shou, B. Qin, T. Liu, D. Jiang, and M. Zhou, ``CodeBERT: A pre-trained model for programming and natural languages,'' in \emph{Findings of EMNLP}, 2020, pp.~1536--1547.
\bibitem{graphcodebert} D. Guo, S. Ren, S. Lu, Z. Feng, D. Tang, S. Liu, L. Zhou, N. Duan, A. Svyatkovskiy, S. Fu et al., ``GraphCodeBERT: Pre-training code representations with data flow,'' in \emph{ICLR}, 2021.
\bibitem{unixcoder} D. Guo, S. Lu, N. Duan, Y. Wang, M. Zhou, and J. Yin, ``UniXcoder: Unified cross-modal pre-training for code representation,'' in \emph{Proc. ACL}, 2022, pp.~7212--7225.
\bibitem{vulbench} Z. Gao, H. Wang, Y. Zhou, W. Zhu, and C. Zhang, ``How far have we gone in vulnerability detection using large language models,'' \emph{arXiv:2311.12420}, 2023.
\bibitem{vulgate} R. Safdar, D. Mateen, S. T. Ali, U. Ashfaq, and W. Hussain, ``Data and context matter: Towards generalizing AI-based software vulnerability detection,'' \emph{arXiv:2508.16625}, 2025.
\bibitem{fish} Y. Shi, J. Seely, P. Torr, N. Siddharth, A. Hannun, N. Usunier, and G. Synnaeve, ``Gradient matching for domain generalization,'' in \emph{ICLR}, 2022.
\bibitem{dgsurvey1} K. Zhou, Z. Liu, Y. Qiao, T. Xiang, and C. C. Loy, ``Domain generalization: A survey,'' \emph{IEEE Trans. Pattern Anal. Mach. Intell.}, vol.~45, no.~4, pp.~4396--4415, 2023.
\bibitem{dgsurvey2} J. Wang, C. Lan, C. Liu, Y. Ouyang, T. Qin, W. Lu, Y. Chen, W. Zeng, and P. S. Yu, ``Generalizing to unseen domains: A survey on domain generalization,'' \emph{IEEE Trans. Knowl. Data Eng.}, vol.~35, no.~8, pp.~8052--8072, 2023.
\end{thebibliography}

\end{document}
